\documentclass[UTF8,a4paper,12pt,fontset=none]{ctexart}

\usepackage{geometry}
\usepackage{array}
\usepackage{longtable}
\usepackage{hyperref}
\usepackage{xcolor}

\setCJKmainfont{Songti SC}
\setCJKsansfont{Songti SC}
\setCJKmonofont{Songti SC}
\setmainfont{Times New Roman}
\setsansfont{Arial}
\setmonofont{Menlo}

\geometry{left=2.5cm,right=2.5cm,top=2.5cm,bottom=2.5cm}
\hypersetup{
  colorlinks=true,
  linkcolor=black,
  urlcolor=blue
}

\title{\textbf{PetFlow 作业报告}}
\author{}
\date{}

\begin{document}

\maketitle

\section{程序功能介绍}

PetFlow 是一个基于任务图的桌面工作流管理程序。与普通待办清单按列表排列任务不同，本项目把任务、习惯、资料、检查点和奖励抽象为图中的节点，把依赖、循环、推荐和触发关系抽象为边。用户可以在画布上直观看到任务之间的先后关系，并通过桌宠和 Agent 辅助完成任务规划、推荐和复盘。

项目使用 Python 3.12 开发，界面基于 Tkinter 和 Canvas 实现，数据通过 JSON 文件进行本地持久化。程序入口为 \texttt{src/petflow/main.py}，运行命令为：

\begin{verbatim}
PYTHONPATH=src python -m petflow.main
\end{verbatim}

程序主要功能如下：

\begin{enumerate}
  \item \textbf{任务图编辑。}用户可以创建、编辑、删除和拖拽节点，也可以创建、编辑、删除节点之间的边。节点支持任务、Routine、资源、检查点、奖励等类型；边支持依赖、Routine、推荐、触发等类型。依赖边会进行成环检测，避免出现逻辑错误。

  \item \textbf{节点信息管理。}每个节点可以记录标题、描述、状态、优先级、预计耗时、实际耗时、标签、检查清单、附件、资源路径、重复规则、下次到期时间等信息。选中节点后，右侧编辑面板可直接修改这些字段。

  \item \textbf{七日议程。}左侧面板展示未来 7 天的任务和 Routine 到期情况，帮助用户把图结构中的任务转化为近期计划。

  \item \textbf{本地推荐。}推荐引擎会根据依赖是否完成、节点状态、优先级、Routine 是否到期、当前焦点任务等因素，推荐下一步最适合执行的节点。

  \item \textbf{桌宠交互。}桌宠不是单纯装饰，而是和任务图状态联动。当用户开始任务、完成任务或请求推荐时，桌宠会移动到对应节点附近，并显示提示气泡。

  \item \textbf{Agent 辅助规划。}程序支持 OpenAI 兼容接口，也支持 mock 模式。用户可以通过自然语言生成任务图，或让 Agent 拆分已有节点。Agent 输出不会直接修改图，而是先转成 proposal，经字段校验和预览后，再通过图服务应用到任务图。

  \item \textbf{系统增强能力。}程序提供剪贴板捕获功能，可将剪贴板中的 URL 或文本转成资源节点；Focus Mode 提供专注计时和前台窗口检测的扩展接口；复盘功能根据历史记录和节点状态生成项目进展总结。

  \item \textbf{保存与加载。}任务图、桌宠状态、工作区状态和历史记录都可保存为 JSON。程序支持默认图文件、样例图和用户保存文件，便于继续编辑和演示。
\end{enumerate}

\section{项目各模块与类设计细节}

项目整体采用分层设计，核心目录为 \texttt{src/petflow/}。主要分为 UI 层、Application 层、Domain 层、Repository 层、Service 层、Agent 层和 System 集成层。这样的划分使界面交互、业务规则、数据存储和外部能力互相隔离，便于多人协作和单元测试。

\subsection{Domain 层：核心数据模型与规则}

Domain 层位于 \texttt{src/petflow/domain/}，不依赖 Tkinter、文件系统或网络请求，是项目中最稳定的核心部分。

\texttt{entities.py} 定义主要实体类：

\begin{itemize}
  \item \texttt{Node}：表示图中的节点。字段包括 \texttt{id}、\texttt{type}、\texttt{title}、\texttt{description}、\texttt{status}、\texttt{priority}、\texttt{estimated\_minutes}、\texttt{actual\_minutes}、坐标、标签、附件、重复规则、资源路径、检查清单等。\texttt{to\_dict()} 和 \texttt{from\_dict()} 负责 JSON 序列化与反序列化。
  \item \texttt{Edge}：表示节点之间的关系。包含 \texttt{source}、\texttt{target}、\texttt{type}、\texttt{label}、\texttt{weight} 等字段。
  \item \texttt{ChecklistItem}：表示节点中的检查项。
  \item \texttt{ProjectMetadata}：保存项目名称、描述、schema 版本和更新时间等元数据。
  \item \texttt{PetState}：保存桌宠当前节点、位置、状态、心情和气泡文本。
  \item \texttt{WorkspaceState}：保存当前选择、缩放、平移、主题和专注模式等界面状态。
\end{itemize}

\texttt{enums.py} 用枚举集中管理字符串常量，例如 \texttt{NodeType}、\texttt{NodeStatus}、\texttt{EdgeType}、\texttt{RepeatType}、\texttt{ResourceType}、\texttt{PetStateType} 和 \texttt{EventType}。这样可以避免在代码中散落 \texttt{"todo"}、\texttt{"dependency"} 等硬编码字符串。

\texttt{graph.py} 中的 \texttt{GraphModel} 是任务图模型，维护所有节点、边、桌宠状态、工作区状态和历史记录。它提供 \texttt{add\_node()}、\texttt{update\_node()}、\texttt{remove\_node()}、\texttt{add\_edge()}、\texttt{update\_edge()}、\texttt{remove\_edge()}、\texttt{predecessors()}、\texttt{successors()} 等方法。图模型会检查节点和边 ID 是否重复、边端点是否存在、删除节点时是否同步删除关联边，以及依赖边是否会造成环。

\texttt{rules.py} 和 \texttt{routine.py} 封装领域规则。依赖关系的成环检测放在领域层，Routine 的下一次到期时间计算也放在领域层，避免 UI 回调中混入业务规则。

\subsection{Application 层：用例编排与事件分发}

Application 层位于 \texttt{src/petflow/app/}，负责把 UI 操作转换为对领域模型的安全修改。

\texttt{GraphService} 是最重要的应用服务，也是 UI 和 Agent 修改任务图的统一入口。它封装了创建节点、创建资源节点、添加附件、更新节点详情、更新节点状态、移动节点、删除节点、创建边、更新边、删除边、设置当前节点等操作。每次修改后，\texttt{GraphService} 会发布事件，便于桌宠、界面刷新和其他服务响应。

\texttt{GraphService.update\_node\_status()} 对状态流转做了集中处理。当节点被标记为 done 时，会写入完成时间；如果是 Routine 节点，还会更新 \texttt{last\_completed\_at}、\texttt{next\_due\_at} 和连续完成次数 \texttt{streak}，并记录历史事件。

\texttt{EventBus} 是同步事件总线。它通过发布和订阅 \texttt{DomainEvent} 解耦模块，例如节点更新后桌宠可以自动响应，而不需要 UI 直接调用桌宠逻辑。

\texttt{AppContext} 是应用上下文，用于集中创建和保存 \texttt{GraphModel}、\texttt{GraphService}、\texttt{StorageService}、\texttt{RecommendationEngine}、\texttt{AgendaService}、\texttt{ReviewService}、\texttt{PetService} 等对象。主窗口只需要持有 \texttt{AppContext}，就可以访问各项服务。

\subsection{Repository 与存储层}

Repository 层位于 \texttt{src/petflow/repositories/}，Service 层中的 \texttt{StorageService} 对其进行封装。

\texttt{GraphRepository} 是抽象接口，定义 \texttt{load()} 和 \texttt{save()} 方法。\texttt{JsonGraphRepository} 是当前实现，负责从 JSON 文件读取任务图，或把 \texttt{GraphModel.to\_dict()} 的结果写入本地文件。默认图文件路径为 \texttt{data/graph.json}。

这种设计避免 UI 直接读写文件。以后如果要支持“另存为”“导入”“导出”或其他存储格式，只需要替换或扩展 Repository 实现。

\subsection{Service 层：业务能力服务}

Service 层位于 \texttt{src/petflow/services/}，用于放置不直接属于实体模型、但又需要被多个界面入口复用的业务逻辑。

\texttt{RecommendationEngine} 负责推荐下一步任务。它先排除已完成和 blocked 节点，再检查依赖前置是否完成。评分时会考虑优先级、doing 状态、paused 状态、Routine 到期、奖励节点和当前焦点节点，最后返回分数最高的候选节点，并能生成推荐理由。

\texttt{AgendaService} 负责生成未来若干天的议程。它读取节点的 \texttt{next\_due\_at} 和重复规则，把任务按日期归组，并按优先级和标题排序。资源节点和已完成节点不会进入议程。

\texttt{RoutineService} 负责 Routine 到期判断和批量刷新。\texttt{ReviewService} 根据图历史、节点状态和 Routine 信息生成复盘数据。\texttt{ResourceService} 整理资源节点打开、复制和预览所需的信息。\texttt{GraphLayoutService} 负责节点自动排列和 Agent 生成子图后的布局调整。

\texttt{PetService} 负责桌宠状态响应。它订阅节点更新事件，当节点进入 doing 状态时移动到该节点，当节点完成时调用推荐引擎寻找下一步，并把桌宠设置为 happy 或 idle 状态。\texttt{PetService} 只更新 \texttt{PetState} 并发布 \texttt{PET\_MOVED} 事件，具体绘制交给 UI 层。

\texttt{MascotService} 负责加载桌宠配置。它从 \texttt{assets/mascots/} 和 \texttt{data/mascots/} 中读取 \texttt{mascot.json}，检查图片资源是否存在，并生成 \texttt{MascotConfig}。

\subsection{UI 层：Tkinter 界面与交互}

UI 层位于 \texttt{src/petflow/ui/}。它的职责是显示数据、收集用户操作，并把操作交给应用服务处理。

\texttt{MainWindow} 是主窗口类，负责创建 Tk 根窗口、顶部工具栏、推荐横幅、左侧议程面板、中心图画布、右侧编辑/伴侣面板、主题切换、Focus Mode、保存加载、Agent 对话框入口等。它还负责在文件加载后重建 \texttt{AppContext}，保证界面和服务使用同一份图模型。

\texttt{GraphCanvas} 继承自 \texttt{tk.Canvas}，负责图的可视化绘制和交互。它维护 canvas item 到 node/edge ID 的映射，支持节点绘制、边绘制、节点拖拽、选择、高亮、缩放、平移、边创建模式、右键菜单、完成动效和图展开动画。Canvas 不直接绕过服务修改业务数据，而是通过 \texttt{GraphService} 完成节点移动、删改和连边。

\texttt{InspectorPanel} 是右侧内联编辑器。选中节点或边后，用户可以在这里修改标题、类型、状态、优先级、预计耗时、实际耗时、描述、标签、资源路径、重复设置、到期日期以及边类型和边标签。它同样通过 \texttt{GraphService} 保存修改，并保留高级对话框作为补充入口。

\texttt{AgendaPanel} 展示未来 7 天计划和保存图文件中的任务。\texttt{PetAssistantPanel} 提供桌宠伴侣聊天和 Agent 辅助入口。\texttt{AgentDialog} 负责展示 Agent 生成的 proposal 预览，并在用户确认后调用 \texttt{AgentExecutor} 应用。\texttt{SettingsDialog} 用于配置 API Key、Base URL、模型和 mock 模式。\texttt{IconButton}、\texttt{components.py}、\texttt{theme.py} 和 \texttt{fonts.py} 提供可复用的界面组件、主题色和字体设置。

\subsection{Agent 层：自然语言规划与校验执行}

Agent 层位于 \texttt{src/petflow/agent/}，分为 Prompt、客户端、proposal 校验和执行四部分。

\texttt{PromptBuilder} 负责根据当前图状态和用户输入构造提示词。\texttt{AgentClient} 负责访问 OpenAI 兼容 API，支持 chat completions 和 responses 两种接口，也支持无 API Key 时回退 mock 模式。它会要求模型返回 JSON，并对响应进行 JSON 解析。

\texttt{AgentProposalValidator} 负责校验 Agent 输出。它限制最大节点数，校验节点标题、类型、状态、优先级、耗时、重复规则、资源类型、坐标、边端点等字段，并兼容部分常见别名，例如 \texttt{completed} 可归一化为 \texttt{done}。

\texttt{AgentExecutor} 负责把校验后的 proposal 应用到图中。它不会直接修改 \texttt{GraphModel}，而是调用 \texttt{GraphService.create\_node()}、\texttt{create\_edge()}、\texttt{update\_node\_detail()} 和 \texttt{delete\_node()} 等方法，因此仍然会经过 ID 生成、字段校验、成环检测和事件发布。对于节点拆分场景，Executor 会自动在父节点和新节点之间添加推荐边。

\subsection{System 集成层}

System 层位于 \texttt{src/petflow/system/}，用于封装和操作系统相关的能力。

\texttt{ClipboardWatcher} 提供一次性剪贴板捕获。它读取剪贴板文本，判断内容是否为 URL，并生成 \texttt{ClipboardCapture}，后续由 \texttt{GraphService.create\_resource\_node()} 转成资源节点。

\texttt{FocusMonitor} 提供前台窗口标题检测。目前该能力只在 Windows 且安装 \texttt{pywin32} 时可用；在 macOS 或未安装依赖时会安全返回不可用，不影响主程序启动。

\section{小组成员分工情况}

\begin{longtable}{|p{3cm}|p{4cm}|p{7cm}|}
\hline
\textbf{成员} & \textbf{主要负责内容} & \textbf{具体工作} \\
\hline
成员 A & 待填写 & 待填写 \\
\hline
成员 B & 待填写 & 待填写 \\
\hline
成员 C & 待填写 & 待填写 \\
\hline
\end{longtable}

\section{项目总结与反思}

本项目的主要收获是把“待办事项管理”设计成了一个可视化任务图系统。相比简单列表，图结构更适合表达前置依赖、循环 Routine、任务拆分和推荐路径。项目中还加入了桌宠、七日议程、Agent 辅助规划、剪贴板资源捕获和复盘等功能，使程序既有可演示性，也有一定实用价值。

从工程实现上看，分层设计带来了明显好处。\texttt{GraphModel} 只关心图结构和领域规则，\texttt{GraphService} 统一处理修改入口，UI 层只负责交互和展示，Agent 层也不能绕过服务直接改数据。这样做虽然前期代码量更多，但后期增加推荐、桌宠、日程和 Agent 时，不需要反复修改底层结构。

项目也暴露出一些不足。第一，UI 层功能较多，\texttt{MainWindow} 和 \texttt{GraphCanvas} 的代码仍然偏大，后续可以继续拆分为更小的控制器或组件。第二，Agent 生成结果虽然已经做了校验，但对复杂自然语言目标的布局和任务粒度控制还有改进空间。第三，Focus Mode 的系统检测能力目前主要面向 Windows，在跨平台体验上还不完整。第四，桌宠交互已经能响应任务状态，但还可以加入更丰富的动画和更细致的反馈策略。

后续如果继续完善，可以优先从三个方向推进：一是继续优化 UI 结构和视觉细节，提高长期使用时的效率；二是增强 Agent proposal 的预览、撤销和冲突处理能力；三是完善跨平台系统能力，让剪贴板、专注检测、文件资源和复盘形成更完整的工作流闭环。

\end{document}
